\newcommand\ndb{qweqwe}
\renewcommand\fgame{\Gcal}
\newcommand\wellbound[1]{\triangleleft_{#1}}

Given an FLTS $L$, a configuration $(s_0,H_0)$ in $L$ and a firm closed formula $\phi_0$,
the \emph{model checking} problem asks whether $(s_0,H_0)\in\sem[]{\phi_0}$. To solve this problem, in this section we introduce parity games that capture satisfiability of formulas in given configurations. To establish this connection it will be useful to consider formulas $\phi_0$ which may not be closed.
We start off with a general definition of parity games in nominal sets.

\begin{definition}\label{def:paritygames}
  A \boldemph{(nominal) parity game} is a tuple $\Gcal=\langle V,V_A,V_D,E,\Omega\rangle$ where $\langle V,E\rangle$ is a directed graph, with vertices called \emph{positions} and edges called \emph{moves}, $V=V_A\uplus V_D$ is a partitioning of positions into \emph{Attacker} and \emph{Defender} ones respectively, and $\Omega:V\to\{0,\dots,d\}$ is a \emph{ranking function} for some maximum rank $d\in\mathbb{N}$.\\
  In addition, $V$ is a finitely supported subset of some nominal set $\hat V$ and, for all $P_1,P_2\in V$ and $P_1',P_2'\in\hat V$, if $(P_1,P_1')\nomeq (P_2,P_2')$ then:
  \begin{itemize}
  \item if $P_1\in V_\chi$ then $P_2\in V_\chi$ (for $\chi\in\{A,D\}$), and $\Omega(P_1)=\Omega(P_2)$;
  \item if $(P_1,P_1')\in E$ then $(P_2,P_2')\in E$.
  \end{itemize}
\end{definition}

Given a game $\Gcal$, we may write $\pos(\Gcal)$ for its set of positions $V$.
Intuitively, the nominal conditions in \cref{def:paritygames} say that the game seen from two nominally equivalent positions is the same up to permutation.
Note that $\Gcal$ is not equivariant in general, and its support is the support of $V$.
%Given a game $\Gcal$ as above, its \emph{nominal closure} is the game $\Cl(\Gcal)=\langle \Cl(V),\Cl(V_A),\Cl(V_D),\Cl(E),\Cl(\Omega)\rangle$ and is equivariant.

Starting from a root position $P_0\in V$, a \emph{play} is a path within $\Gcal$, i.e.\ a finite or infinite sequence $\Pi=P_0,P_1,\dots$ such that $(P_i,P_{i+1})\in E$ for each $i$. A play is \emph{complete} if it ends in a leaf. Defender (Attacker) wins a complete play if the last position in it belongs to Attacker (resp.~Defender). Defender (Attacker) wins an infinite play if the highest position rank in it is even (resp.~odd). We say that  $P_0$ is a \emph{winning position} for Defender (or Defender \emph{wins from} $P_0$) if there is a partial function $\Theta: V_D\rightharpoonup E$ (called a \emph{strategy}) such that Defender wins  every complete or infinite play $P_0,P_1,\dots$ such that for each $P_i\in V_D$ we have $\Theta(P_i)=(P_i,P_{i+1})$.

The development of parity games for \fhml\ follows
parity games for the $\mu$-calculus, cf.~\cite{EmersonJutla91}. We start by
considering the negation-free fragment of our logic. In the sequel, we shall consider the derived connectives of \cref{rem:derived} as primitives (and exclude negation).

\begin{definition}
  A formula is called \emph{negation-free}
  if it contains no negations. For each formula $\phi$, we
write $!\phi$ for its negation-free variant. This can be constructed by induction using standard duality rules for pushing negation inside boolean, modal and recursion constructors, along with the rules: \
$ !(\lnot\lnot\phi)={!\phi},\
        !(\lnot \fresh{x} \phi) = \fresh{x} !(\lnot \phi) 
$.
\end{definition}


\begin{restatable}{lemma}{negfreesize}\label{lem:negfreesize}
    For any formula $\phi$, $|!(\phi)|$ is bounded by $|\phi| - n$, where $n$ is the number of negations in $\phi$. 
\end{restatable}

For the remainder of this section, and unless specified otherwise, any \fhml\ formulas used are going to be assumed to be firm and negation-free.

% In our definitions it will be useful to use substitutions of recursion formulas for recursion variables. As the latter have arities, such substitutions are parametric on the arguments variables are applied on. To denote substitutions of this form, we use the following notation:
% \[
%     \{\sigma X(\vec x).\psi / X \} = \{(\sigma X(\vec x).\psi)(\vec a) / X(\vec a) \mid \vec a \in \Atoms^{\ar(X)}\}
% \]
% with $\sigma\in\{\mu,\nu\}$.
We shall make use of the (recursion) \emph{alternation depth} of formulas. At this point it is important that the formulas we examine have a unique specification for each of their bound variables. More specifically, for any formula $\phi$ we stipulate that:
\begin{itemize}
\item the bound and free recursion variables of $\phi$ be disjoint;
\item if $\sigma_1 X(\vec x_1).\phi_1,\sigma_2 X(\vec x_2).\phi_2$ are subformulas of $\phi$ then $(\sigma_1,\vec x_1,\phi_1)=(\sigma_2,\vec x_2,\phi_2)$.
\end{itemize}
Note that these conditions are not unusual (cf.~\cite{KlinL17}). They are not restricting our logic either as any formula has an $\alpha$-equivalent one in the form above. At the same time, they allow us to match bound recursion variables with their binding subformulas.

For value variables, on the other hand, we simply disallow a binder on a variable $x$ to be within the scope of another binder on $x$.

\begin{definition}[\cite{BradfieldW18, EmersonJS01}]\label{def:altDepthNom}
    Let $\phi$ be a \fhml\ formula. 
    We define the \emph{dependency order} on bound variables of $\phi$ as the smallest partial order $\leq_\phi$ such that\
    $X \leq_\phi Y$\
    if $X(\vec u)$ occurs free in some $\phi'$ such that $\sigma Y(\vec x) .\phi'$ appears in $\phi$ (for some $\vec u$). 
    \\
    The \boldemph{alternation depth} of a 
    bound variable $X$ in $\phi$ is defined as the maximal length of a chain $X=X_1 \leq_\phi \dots \leq_\phi X_n$ such that:
    \begin{itemize}
        \item if $X$ is a $\mu$-variable, each $X_i$ is a $\mu$-variable if $i$ is odd, and a $\nu$-variable otherwise; 
    \item if $X$ is a $\nu$-variable,  each $X_i$ is a $\nu$-variables if $i$ is odd, and a $\mu$-variable otherwise. 
    \end{itemize}
    The {alternation depth} of a formula $\phi$ (denoted \emph{adepth($\phi$)}) is the maximum alternation depth of variables bound in $\phi$, or zero if there are no fixpoints.
\end{definition}

Parity games for \fhml\ have positions of the form $(s, H, \xi, \psi)$, with:
\begin{itemize}
    \item $(s, H) \in \fstates$ being the current configuration that is examined;
    \item $\xi$ being a $\Ucal$-variable assignment;
    \item $\psi$ being the current formula. 
\end{itemize}
We define games with respect to a \emph{root position} corresponding to the model-checking problem we aim to decide. The positions and moves are given inductively using \emph{game rules}.

\begin{definition}\label{def:parityG}
    Let $\phi_0$ be a firm negation-free \fhml\
    formula,
    $\xi$ a $\Ucal$-variable assignment containing the free recursion variables of $\phi_0$, $\Lcal$ an FLTS and $(s_0,H_0)$ a configuration in $\Lcal$.
    The \boldemph{game} $\fgame(\Lcal, \phi_0, \xi, s_0, H_0)$ has root position $(s_0, H_0, \xi, \phi_0)$ and game rules:
    \begin{itemize}
        \item $(s,H,\xi, a = a)$ belongs to Attacker, whereas
 $(s,H,\xi, a = b)$ to Defender;
        \item $(s,H, a \neq a)$ belongs to Defender, whereas 
 $(s,H, a \neq b)$  to Attacker;
        \item $(s,H,\xi, X(\vec a))$ belongs to Attacker if $(s, H)\in \xi(X)(\vec a)$, and to Defender otherwise;
        \item $(s,H,\xi, \modal{\ell}\phi) \rightarrow (s',H',\xi, \phi)$, if $(s, H) \xrightarrow{\ell} (s', H')$, belongs to Defender if $\modal{\ell}=\diam{\ell}$, and to Attacker if $\modal{\ell}=[\ell]$;
        \item $(s,H,\xi, \phi_1 \odot \phi_2) \rightarrow (s,H,\xi,\phi_i)$, $i \in \{1,2\}$, belongs to Defender if $\odot=\lor$, and to Attacker if $\odot=\land$;
        % \item $(s,H,\xi, \phi_1 \wedge \phi_2) \rightarrow (s,H,\xi, \phi_i)$, $i \in \{1,2\}$, played by Attacker,
        \item $(s,H,\xi, \bigodot\nolimits_x \phi) \rightarrow (s,H,\xi,\phi\{a/x\})$, $a \in \Atoms$, belongs to Defender if $\bigodot=\bigvee$, and to Attacker if $\bigodot=\bigwedge$;
        % \item $(s,H,\xi, \bigwedge\nolimits_x \phi) \rightarrow (s,H,\xi,\phi\{a/x\})$, $a \in \Atoms$, played by Attacker,
        % \item $(s,H,\xi, [\ell]\phi) \rightarrow (s',H',\xi, \phi)$, $\forall (s', H'). (s, H) \xrightarrow{\ell} (s', H')$, played by Attacker,
        \item $(s,H,\xi,\fresh{x}\phi) \rightarrow (s, H, \xi, \phi\{a/x\}),\ a \# \phi,\xi,H$, belongs to either;
        \item $(s,H,\xi, (\sigma X(\vec x).\phi)(\vec a)) \rightarrow (s, H, \xi,(\phi\{\sigma X(\vec x).\phi / X\})\{\vec a / \vec x\})$, $\sigma\in\{\mu,\nu\}$, belongs to either.
        % \item $(s,H,\xi, (\nu X(\vec x).\phi')(\vec a)) \rightarrow (s, H, \xi, (\phi'\{\nu X(\vec x).\phi' / X\})\{\vec a / \vec x\})$, played by either.
    \end{itemize}
    % In the special cases where $\phi$ is a closed formula (i.e., with no free recursive variables $X$), we suppress $\xi$ everywhere as it is not used and simply write $\fgame(\Lcal, \phi, s, H)$,
  % For winning conditions we stipulate that, in a finite play, a player wins if the other player is required but unable to make a move.
  % In an infinite play, the player who owns the highest infinitely recurring rank wins, where Defender owns even-ranked positions, and Attacker owns odd-ranked ones.
      The rank of a position with formula $\phi$ is given by $\Omega(\phi)$:
    \[        \Omega(\phi) =
      \begin{cases} 2 \cdot \lfloor adepth(X)/2 \rfloor & \text{if $\phi$ is of the form }(\nu X(\vec x).\phi')(\vec a)\\
        2 \cdot \lfloor adepth(X)/2 \rfloor + 1 & \text{if $\phi$ is of the form }(\mu X(\vec x).\phi')(\vec a)\\
         0 &\text{otherwise}
      \end{cases}
    \]
In the case when $\phi_0$ is closed, we denote the game simply by $\fgame(\Lcal, \phi_0, s_0, H_0)$.    
\end{definition}

Note that the formulas appearing inside positions are firm and negation-free (as is $\phi_0$ by assumption). 

%
%  \begin{example}
%One or two examples here
%  \end{example}

Later in this section it will be important to count the orbits of a history-bounded version of the parity game of \cref{def:parityG}. To that effect, it is useful to the formulas appearing in the game, in terms of subformulas of $\phi_0$ and substitutions. This is achieved by the following notion of formula closure. 
  
\begin{definition}\label{def:closFresh}
    Given a \fhml\ formula $\phi_0$, we define its closure $\clos(\phi_0)$ to be the smallest set $\chi$ of triples $(\phi,\gamma,\theta)$ such that $(\phi_0,\varepsilon,\varepsilon) \in \chi$ and:
    \begin{itemize}
        \item if $(\phi_1 \odot \phi_2, \gamma,\theta) \in \chi$, with $\odot\in\{\lor,\land\}$, then $(\phi_1, \gamma,\theta), (\phi_2, \gamma,\theta) \in \chi$;
        \item if $(\modal{\ell}\phi, \gamma,\theta) \in \chi$, with $\modal{l}\in\{\diam{\ell},\sq{\ell}\}$, then $(\phi, \gamma,\theta) \in \chi$;
        \item if $(\bigodot_x\phi, \gamma,\theta) \in \chi$, with $\bigodot\in\{\bigvee,\bigwedge\}$, %or $(\fresh{x} \phi', \theta) \in \chi$
          then $(\phi,(\gamma,\{a/x\}),\theta) \in \chi$, for all $a\in\Atoms$;
        \item if $(\fresh{x}\phi,\gamma, \theta) \in \chi$
          then $(\phi,(\gamma,\{a/x\}), \theta) \in \chi$, for all $a\#\phi\{\theta\}\{\gamma\}$;
        \item if $((\sigma X(\vec x).\phi)(\vec u),\gamma, \theta) \in \chi$, with $\sigma\in\{\mu,\nu\}$, then $(\phi,(\gamma,\{\vec c/\vec x\}), (\theta,\{\sigma X(\vec x).\phi / X\}))\in\chi$, for all $\vec c \in \Atoms^{\ar(X)}$.
        % \item if $((\nu X(\vec x).\phi')(\vec a), \theta) \in \chi$ then $\{(\phi'\{\vec c/\vec x\}, \theta[(\nu X(\vec x).\phi')\{\theta\} / X]) \mid \vec c \in \Atoms^{\ar(X)}\} \subseteq \chi$. 
    \end{itemize}
    We write $\clos^*(\phi_0) = \{\phi\{\theta\}\{\gamma\}\mid (\phi,\gamma, \theta) \in \clos(\phi_0)\}$.
\end{definition}

The following lemma shows that the formulas appearing in the game $\fgame(\Lcal,\phi_0, s_0,H_0)$ can be traced back to the closure. It moreover provides some bounds on their orbits and supports. 

\begin{restatable}{lemma}{freshClosProps}\label{lem:freshClosProps}
Given a closed formula $\phi_0$, an FLTS $\Lcal$ and configuration $(s_0,H_0)$:
    \begin{enumerate}
        % \item\label{lem:freshClosProps:1} for all $(\phi,\gamma,\theta)\in\clos(\phi_0)$, the free variables of $\phi$ are included in $\dom(\theta)\uplus \cal Z$ and all formulas in $\rng(\theta)$ are fixpoints and have free variables in $\dom(\theta)\uplus\Zcal$;
        \item\label{lem:freshClosProps:2} for all positions $(s,H,\phi)$ in $\fgame(\Lcal,\phi_0, s_0,H_0)$ we have that $\phi\in\clos^*(\phi_0)$; 
        \item\label{lem:freshClosProps:3} $|\orb(\clos^*(\phi_0))|$ is bounded by $|\phi_0|\cdot \frac{(|\supp(\phi_0)| + \|\phi_0\|)!}{|\supp(\phi_0)|!}$;
\item\label{lem:freshClosProps:4} for all $(\phi,\gamma,\theta)\in\clos(\phi)$, $|\supp(\phi,\gamma,\theta)| \le |\supp(\phi_0)| + \|\phi_0\|$. 
%\ntnote{this needs to be proven still}
    \end{enumerate}
\end{restatable}


% \begin{definition}
%     Let $\phi$ be a \fhml\ formula, $\Lcal$ be an F-LTS and $\S = |\supp(\phi)| + \|\phi\| + \max(\{(|\supp(s')|) \mid s' \in \Scal\})$.
%     For any $H \subseteq \Atoms$ and $s \in \Scal$, let us denote $H \upharpoonright (s, \phi, \S) = H$ if $|H| \le \S$, otherwise:
%     \[
%         H \upharpoonright (s, \phi, \S) = H' \text{ such that } \supp(\phi, s) \cap H \subseteq H' \subseteq H \text{ and } |H'| = \S + 1
%     \]
%     Given a game $\Gcal = \Gcal(\Lcal, \phi, \xi, s)$, we construct the game $\Gcal_\S$ from $\Gcal$ as follows:
%     \begin{itemize}
%         \item $\pos(\Gcal_\S) = \{\orb[\supp(\phi, s)](s, H\upharpoonright (s, \phi, \S), \phi) \mid (s, H, \phi) \in \pos(\Gcal)\}$
%         \item for every move $(s, H, \phi) \rightarrow (s', H', \phi')$, then there exists a move $\orb[\supp(\phi, s)](s, H \upharpoonright (s, \phi, \S), \phi) \rightarrow \orb[\supp(\phi', s')](s', H' \upharpoonright (s, \phi, \S), \phi')$
%     \end{itemize}
%     The root of $\Gcal_\S$ is the same as that of $\Gcal$.
% \end{definition}

% \begin{definition}
%     Let $\phi$ be a \fhml\ formula, $\Lcal$ be an F-LTS and $\S = |\supp(\phi)| + \|\phi\| + \max(\{(|\supp(s')|) \mid s' \in \Scal\})$.
%     Given a game $\Gcal = \Gcal(\Lcal, \phi, \_, s)$, we construct the game $\Gcal_\S = \Gcal_\S(\Lcal, \phi, \_, s)$ as follows:
%     \begin{itemize}
%         \item $\pos(\Gcal_\S) = \{(s, H, \phi) \in \pos(\Gcal) \mid |H| \le \S\} \cup \{(s, H', \phi) \mid (s, H, \phi) \in \pos(\Gcal) \land |H| > \S \land \supp(\phi, s) \cap H \subseteq H' \subseteq H \land |H'| = \S + 1\}$
%     \end{itemize}
%     For every move $(s_1, H_1, \phi_1$
%     The root of $\Gcal_\S$ is the same as that of $\Gcal$.
% \end{definition}

We proceed to showing the correspondence between the semantics of an \fhml\ formula and its corresponding parity game.
\begin{restatable}{theorem}{FnomSatPar}\label{thm:FnomSatPar} %THEOREM 6 - BRADFIELDW18
    Given a \fhml\ formula $\phi_0$,
    $\Ucal$-variable assignment $\xi$,
    FLTS $\Lcal$ and $(s_0, H_0)$ from $\Lcal$, we have $(s_0, H_0) \in \sem{\phi_0}$ iff Defender wins from $(s_0, H_0,\xi, \phi_0)$ in $\fgame(\Lcal, \phi_0, \xi, s_0,H_0)$.
\end{restatable}

We have thus reduced model checking of formulas to winning in parity games. 
At this stage, we can restrict our attention to closed (firm negation-free) formulas $\phi_0$.
The next step is deciding
winning regions of our parity games. The latter requires the games to be finite, whereas ours are not so. We therefore reduce our parity games to equivalent finite ones.
For this to be possible, we restrict our attention to FLTSs that are fresh-orbit-finite.

For brevity, in the remainder of this section we shall say that $(\Lcal,\phi_0,s_0,H_0)$ is a \boldemph{setup of grade $N$} if:
\begin{itemize}
\item $\Lcal$ is a fresh-orbit-finite FLTS with set of states $\Scal$, and $(s_0,H_0)\in\hat \Scal$;
\item $\phi_0$ is a closed (firm negation-free) \fhml\ formula;
\item $N = |\supp(\phi_0)| + \|\phi_0\| + \regindex(\Scal)$. 
\end{itemize}
%
Given a setup $(\Lcal,\phi_0,s_0,H_0)$ of grade $N$, we first present a version of $\fgame(\Lcal,\phi_0,s_0,H_0)$ with bounded histories.\ntnote{$N$ is used for sets of names earlier, can be confusing.} 
The idea is to restrict the size of the history to the maximum number of names that are available to the formula and the states of the FLTS, and one extra to represent a fresh name. This amounts to a bound of $N+1$ names. 
When a new name is added and the size of the history is greater than $N$, then names that are in the history but no longer present in the formula are forgotten. 
Each position in this game is a tuple $(s, H, \phi)$ with $|H|\leq N+1$.
% \begin{itemize}
%     \item $(s, H) \in \fstates$ being the current configuration being examined and $\phi$ a \fhml\ formula;
%     \item $|H| \le \S + 1$ and $\supp(s) \subseteq H$.
% \end{itemize}
\begin{definition}
Let $(\Lcal,\phi_0,s_0,H_0)$ be a setup of grade $N$. 
The \boldemph{history-bounded game} $\Gcal_\S(\Lcal,\phi_0, s_0, H_0)$ has root
 position $(s_0, H_0', \phi_0)$, for some $H_0'\wellbound{N}(s_0,\phi_0,H_0)$, and game rules as in \cref{def:parityG}, apart from the rule for modal connectives:
    \begin{itemize}
 %        \item $(s,H, a = a)$ belongs to Attacker, whereas 
 % $(s,H, a = b)$ belongs to Defender;
 %        \item $(s,H, a \neq a)$ belongs to Defender, whereas 
 % $(s,H, a \neq b)$ belongs to Attacker;
        \item $(s,H, \modal{\ell}\phi) \rightarrow (s',H', \phi)$, if $\ell=(t, \vec a)$
        and $(s,H) \xrightarrow{t, \vec a} (s',H\cup\{\vec a\})$ with $H'\wellbound{N}(s',\phi,H\cup\{\vec a\})$,
        played by Defender if $\modal{\ell}=\diam{\ell}$, and by Attacker if $\modal{\ell}=\sq{\ell}$;
        \end{itemize}
% \item $(s,H, \phi_1 \odot \phi_2) \rightarrow (s,H,\phi_i)$, $i \in \{1,2\}$, belongs to Defender if $\odot=\lor$, and by Attacker if $\odot=\land$;
%         \item $(s,H, \bigodot\nolimits_x \phi) \rightarrow (s,H,\phi\{a/x\})$, $a \in \Atoms$, played by Defender if $\bigodot=\bigvee$, and by Attacker if $\bigodot=\bigwedge$;
%         \item $(s,H,\fresh{x}\phi) \rightarrow (s, H, \phi\{a/x\}),\ a \#\phi,\xi,H$, played by either;
%         \item $(s,H, (\sigma X(\vec x).\phi)(\vec a)) \rightarrow (s, H,\phi\{\sigma X(\vec x).\phi / X\}\{\vec a / \vec x\})$, $\sigma\in\{\mu,\nu\}$, played by either.
while the ranking function is the same as in \cref{def:parityG}.
The \emph{well-bounding} relation $\triangleleft$ is given by $H'\wellbound{N}(s,\phi,H)$ if $H'\subseteq H$ is such that:
\begin{itemize}
            \item if $|H| \le N$ then $H' = H$,
            \item else $H \cap \supp(s,\phi) \subseteq H' \subseteq H$ and $|H'| = N + 1$.
        \end{itemize}
\end{definition}

We note that, though the positions in the unbounded game have their histories trimmed at modal operators, they remain FLTS configurations, i.e.\ for each $(s,H,\phi)$ we have $\supp(s)\subseteq H$. 
%
We next show that we can decide winning a parity game $\Gcal$ by examining its equivalent history-bounded game $\Gcal_N$.
To do this, we are going to build a bisimulation-like relation $\Rcal$ between $\pos(\Gcal)$ and $\pos(\Gcal_N)$. Since the positions in the two games can have name-mismatches, but essentially do ``the same'' transitions, it is useful to introduce a notion of {bisimilation} up to renaming. Note use of relation composition: $R_1;R_2=\{(x,z)\mid \exists y.\,(x,y)\in R_1\land (y,z)\in R_2\}$.

\begin{definition}
Given parity games $\Gcal_1, \Gcal_2$, we call a relation $R \subseteq \pos(\Gcal_1) \times \pos(\Gcal_2)$ a \emph{bisimulation} if for all $(P_1, P_2) \in R$:
    \begin{itemize}
        \item $P_1$ has the same rank and belongs to the same player as $P_2$;
        \item for each $P_1 \rightarrow P_1'$, there exist some $P_2 \rightarrow P_2'$  such that $(P_1',P_2')\in R$;
        \item for each $P_2 \rightarrow P_2'$, there exist some $P_1 \rightarrow P_1'$ such that $(P_1', P_2') \in R$.
    \end{itemize}
    We say that $P_1$ and $P_2$ are \emph{bisimilar} (denoted ${P_1}\bisim{P_2}$)
    if there is a bisimulation $R$ such that $(P_1, P_2) \in R$. \\
    A \emph{bisimulation up to nominal equivalence (nom-bisimulation)} is defined to be a relation $R$ with the same properties as above, apart from the requirement in the last two bullet points being instead that $(P_1',P_2')\in {\nomeq};R;{\nomeq}$.
    We say that $P_1$ and $P_2$ are \emph{nom-bisimilar} (denoted ${P_1}\nombisim{P_2}$)
    if there is a nom-bisimulation $R$ such that $(P_1, P_2) \in R$. 
\end{definition}

In the following two lemmas, we let $P_1,P_2$ be positions in games $\Gcal_1,\Gcal_2$ respectively.

\begin{lemma}\label{lem:nombisimbisim}
If $P_1\nombisim P_2$ then $P_1\bisim P_2$. %(seen as positions in $\Cl(\Gcal_1),\Cl(\Gcal_2)$ respectively). 
\end{lemma}
\begin{proof}
  Let $R\subseteq\pos(\Gcal_1)\times\pos(\Gcal_2)$ be 
  a nom-bisimulation. We show that $R'={\nomeq;R;\nomeq}\subseteq\pos(\Gcal_1)\times\pos(\Gcal_2)$ is a bisimulation.
Suppose $(P_1,P_2)\in R'$, i.e.\ $\pi_i\cdot P_i=Q_i$ ($i=1,2$) and $(Q_1,Q_2)\in R$, and let $P_1\to P_1'$. We have $P_1,Q_1\in\pos(\Gcal_1)$ and, taking $Q_1'=\pi_1\cdot P_1'$, $(P_1,P_1')\nomeq(Q_1,Q_1')$, hence $Q_1\to Q_1'$ (and $Q_1'\in\pos(\Gcal_1)$). Then, by nom-bisimulation, $Q_2\to Q_2'$ and $(Q_1',Q_2')\in R'$. As before, setting $P_2'=\pi_2^{-1}\cdot Q_2'$, we obtain $P_2\to P_2'$ (and $P_2'\in\pos(\Gcal_2)$). We observe that $(P_1',P_2')\in{\nomeq;R';\nomeq}=R'$.
% Note that nominal closure may take us outside the positions of $\Gcal_i$ (for $i=1,2$), hence the need to consider $\Cl(\Gcal_i)$.
\end{proof}

\begin{restatable}{lemma}{nombisimequiv}\label{lem:nombisimequiv}
If $P_1\bisim P_2$ then Defender wins from $P_1$ iff they win from $P_2$.
\end{restatable}

% Proof attempt, gets overly long
%
% \begin{proof}
% By symmetry, it suffices to show the left-to-right direction. Suppose $(P_1,P_2)\in R$ for some bisimulation $R$ and let $\Theta_1$ be a winning strategy for Defender from $P_1$.
% Let $\iota:\pos(\Gcal_2)\rightharpoonup\pos(\Gcal_1)$ be a function selecting for each $P_2'$ a representative from $R^{-1}(P_2')$ (if $R^{-1}(P_2')\neq\emptyset$).  
% We build the winning strategy $\Theta_2$ from $P_2$ as follows:
% \[
%   \Theta_2(P_2') = \begin{cases}
%                      P_2'' & \text{if $P_1'=\iota(P_2')$ and $\Theta_1(P_1')=P_1\to P_1''$ is matched by $P_2'\to P_2''$ in $R$}\\
%                      \text{undefined} & \text{otherwise}
%                    \end{cases}
% \]
% where $P_2'$ ranges over Defender positions in $\Gcal_2$ and choice is used in the first case above.
% Let us now consider a complete or infinite play $\Pi_2=P_{2,0},P_{2,1},\dots$ in $\Gcal_2$ with $P_{2,0}=P_2$ and following $\Theta_2$. We build a sequence $\Pi_1=P_{1,0},P_{1,1},\dots$ in $\Gcal_1$ inductively by setting $P_{1,0}=P_1$ and, for each $0\leq i<|\Pi_2|-1$, assuming $(P_{1,i},P_{2,i})\in R$:
%     \begin{itemize}
%     \item if $P_{1,i}$ is an Attacker position then we take $P_{1,i+1}$ such that $P_{1,i}\to P_{1,i+1}$ matches $P_{2,i}\to P_{2,i+1}$ in $R$; 
%         \item if $P_{1,i}$ is a Defender position then we take $P_{1,i+1}$ such that $\iota(P_{2,i})\to P_{1,i+1}$ is the transition matched by $P_{2,i}\to P_{2,i+1}$ in the definition of $\Theta_2$;
%         \end{itemize}
%         thus, in both cases, $(P_{1,i+1},P_{2,i+1})\in R$. We note that $\Pi_1$ may not be a play but, for each $i$,
% $(P_{1,i},P_{2,i})\in R$ and 
%         there is $\hat P_{i}\in\pos(\Gcal_1)$ such that $P_{1,i}\bisim\hat P_{i}\to P_{1,i+1}$. Indeed, if $P_{1,i}$ is an Attacker position we simply take $\hat P_{i}=P_{1,i}$, and if $P_{1,i}$ is a Defender position we let $\hat P_{i}=\iota(P_{2,i})$. We can therefore build a play $\Pi_1'=P_{1,0}',P_{1,1}',\dots$ in $\Gcal_1$ inductively by setting $P'_{1,0}=P_{1,0}$ and, for each $0\leq i<|\Pi_1|-1$, assuming $P_{1,i}'\bisim P_{1,i}$ (and, hence, $P_{1,i}'\bisim \hat P_{i}$), pick $P_{1,i+1}'$ such that $P_{1,i}'\to P_{1,i+1}'$ matches $\hat P_i\to P_{1,i+1}$. Thus, for each $i$, $P_{1,i}'\bisim P_{1,i}$ and therefore $P_{1,i}'\bisim P_{2,i}$.
% \end{proof}

We can relate games with their history-bounded counterparts by nom-bisimulations.


\begin{definition}\label{def:theFreshRelation}
Given a setup $(\Lcal,\phi_0,s_0,H_0)$ of grade $N$, consider $\fgame(\Lcal, \phi_0,s_0, H_0)$  and $\Gcal_\S(\Lcal, \phi_0, s_0, H_0)$.
We define the \emph{history-bounding relation} $\Rcal_{(\Lcal, \phi_0, s_0, H_0)} \subseteq \pos(\fgame) \times \pos(\Gcal_\S)$, by setting $((s_1, H_1, \phi_1), (s_2, H_2, \phi_2)) \in \Rcal$ if $s_1 = s_2$, $\phi_1 = \phi_2$ and $H_2\wellbound{N}(s_1,\phi_1,H_1)$.
    % \begin{itemize}
    %     \item if $|H_1| \le \S$ then $H_1 = H_2$,
    %     \item if $|H_1| > \S$ then $\supp(\phi_1) \cap H_1 \subseteq H_2 \subseteq H_1$ and $|H_2| = \S + 1$.
    % \end{itemize}
\end{definition}

Thus, a position $(s,H,\phi)$ in the unbounded game is related to itself if its history $H$ is small, i.e.\ it has no more than $N$ names. If $|H|$ is larger, then we relate $(s,H,\phi)$ to each $(s,H',\phi)$ where $H'\subseteq H$ has $N+1$ names and in particular contains $\supp(s)$ (as $(s,H')$ is a configuration) and all names from $\supp(\phi)$ that appear in $H$. 



\begin{restatable}{lemma}{boundingisnombisim}\label{lem:boundingisnombisim}
    % Let $\hat\phi$ be a \fhml\ formula, $\Lcal$ be an F-LTS, $(\hat s,\hat H)$ a state in $\Lcal$ and $\S = |\supp(\hat \phi)| + \|\hat \phi\| + \regindex(\fstates)$. Pick $\hat H'$ such that 
    % \begin{itemize}
    %   \item if $|\hat H| \le \S$ then $\hat H' = \hat H$,
    %     \item if $|\hat H| > \S$ then $\supp(\hat\phi, \hat s) \cap \hat H \subseteq \hat H' \subseteq \hat H$ and $|\hat H'| = \S + 1$.
    %         \end{itemize}
    %     Let $\fgame = \fgame(\Lcal, \hat \phi, \hat s, \hat H)$ and $\Gcal_\S = \Gcal_\S(\Lcal, \hat \phi, \hat s, \hat H')$, and let $\Rcal_{(\Lcal, \hat \phi, \hat s, \hat H)} \subseteq \pos(\fgame) \times \pos(\Gcal_\S)$ be a history bounded relation.
    % Then, $((\hat s, \hat H, \hat \phi), (\hat s, \hat H', \hat \phi)) \in \Rcal_{(\Lcal, \hat \phi, \hat s, \hat H)}$ and 
    % $\Rcal_{(\Lcal, \hat \phi, \hat s, \hat H)}$ is a bisimulation up to nominal equivalence.
$\Rcal_{(\Lcal, \phi_0, s_0, H_0)}$ is a nom-bisimulation.
\end{restatable}

Combining \cref{thm:FnomSatPar} with \cref{lem:boundingisnombisim,lem:nombisimbisim,lem:nombisimequiv}, we can reduce \fhml\ model-checking to deciding winning a history-bounded parity game. However,
although the size of histories, and therefore of positions, is now bounded, the number of positions can still be infinite due to the names appearing in them. We can resolve this by grouping together positions that are ``equivalent'' up to permuting their names. Formally, this is done by considering games in which positions are orbits of ordinary positions.

% If we have a parity game $\fgame$, with a corresponding history bounded game $\Gcal_\S$, 
% then for the history bounded relation $\Rcal_{(\Lcal, \phi, s, H)} \subseteq \pos(\fgame) \times \pos(\Gcal_\S)$, we show that $\Rcal' = \{(\orbs(P_1), \orbs(P_2)) \mid (P_1, P_2) \in \Rcal_{(\Lcal, \phi, s, H)}\}$ is a bisimulation.

\begin{definition} 
    Given a parity game $\Gcal=\langle V,V_A,V_D,E,\Omega\rangle$, 
    we construct the \boldemph{orbits game} $\gorb(\Gcal)=\langle V',V_A',V_D',E',\Omega'\rangle$ as follows:
    \begin{itemize}
    \item $V' = \{\orbs(P) \mid P \in V\}$ and $V_\chi'=\{\orbs(P) \mid P \in V_\chi\}$ (for $\chi\in\{A,D\}$);
      \item $E'=\{(\orbs(P),\orbs(P'))\mid (P,P')\in E\}$;
        \item  for each $P\in V$, $\Omega'(\orbs(P)) = \Omega(P)$. 
    \end{itemize}
\end{definition}

We can see that $\gorb(\Gcal)$ is well defined. In particular, $\Omega'$ is well defined due to the fact that $\Omega(P_1)=\Omega(P_2)$ whenever $P_1\nomeq P_2$.  
Moreover, we can show that the relation\ $R
=        \{(P, \orbs(P)) \mid P \in \pos(\Gcal)\}$\ is a bisimulation.


\begin{lemma}\label{lem:parityBisim}
  Given $\Gcal$ as above and $P\in\pos(\Gcal)$, $P\bisim\orbs(P)$.
\end{lemma}

Taking stock of the game transformations defined above and the bisimulation relations between them, we can decide winning in a setup of finite grade.
%For the remainder of this section, we fix a firm closed \fhml\ formula $\phi_0$, a fresh-orbit-finite FLTS $\Lcal$ and a configuration $(s_0,H_0)$ of $\Lcal$.

\begin{corollary}\label{cor:winningNom}
Given a setup $(\Lcal,\phi_0,s_0,H_0)$ of grade $N$, Defender wins in $\Gcal(\Lcal,\phi_0,s_0,H_0)$ (from $(s_0,H_0,\phi_0)$) iff they win in $\gorb(\Gcal_N(\Lcal,\phi_0,s_0,H_0))$ from $\orbs(s_0,H_0',\phi_0)$. Hence, it is decidable whether $(s_0,H_0)\in\sem[]{\phi_0}$.
\end{corollary}
\begin{proof}
The first part follows from \cref{thm:FnomSatPar,lem:boundingisnombisim,lem:nombisimbisim,lem:nombisimequiv,lem:parityBisim}. The second part follows from \cref{thm:FnomSatPar} and the fact that 
$\gorb(\Gcal_N(\Lcal,\phi_0,s_0,H_0))$ is a finite game.
\end{proof}

We next calculate upper bounds for the size of $\gorb(\Gcal_N(\Lcal,\phi_0,s_0,H_0))$ and the time needed to construct it. This allows us to give a complexity bound for model checking.
Recall that we set $N=|\supp(\phi_0)|+\|\phi_0\|+\regindex(\Scal)$. In the calculations below we prefer to distinguish between complexity due to $\phi_0$ and $\Lcal$, so instead of $N$ we shall use the measure: $M(\phi_0)=|\supp(\phi_0)|+\|\phi_0\|$. This can be seen as the \emph{nominal potential} of $\phi_0$.%\ntnote{bring in the definition of $\Phi$}

In constructing the orbits game, we need to be able to check if two positions of $\Gcal_N$ are nominally equivalent. The way this is performed is described in \cref{app:orbitGameSize}: given $(s_1,H_1,\phi_1),(s_2,H_2,\phi_2)$, we try to build a permutation $\pi$ such that $\pi\cdot\phi_1=\phi_2$ and, if successful, we extend $\pi$ to some $\pi'$ so that, in addition, $\pi'\cdot s_1=s_2$. The latter part requires some oracle that is able to answer such questions for the given nominal set of states $\Scal$.
%
We let a \emph{permutation oracle} $\Psi_\Lcal$ to be a function that,
given  $s, s' \in \Scal$
	 and  $\vec a, \vec b \in \Atoms^k$ (for some $k$, and each containing distinct names), the oracle returns 
	a permutation $\pi'$ such that $\pi' \cdot s =s'$ and $\pi'\cdot \vec a = \vec b$,
	or $\NO$ if no such permutation exists. 
	We stipulate that the time complexity of $\Psi_\Lcal(\vec a, \vec b, s, s')$ is bounded by some $\Phi_\Lcal(k)$.
%\end{definition}

% \begin{restatable}{lemma}{sizebound}\label{lem:sizebound}
% 	For every \fhml\ formula $\phi$ such that for every subformula $(\mu X(\vec x).\phi')(\vec u)$ of $\phi$ we have $X$ is free in $\phi'$, then $|\supp(\phi)|+2\|\phi\|+\zeta(\phi)\leq|\phi|$ where $\zeta(\phi)$ is the number of free value variables (of the form $x, y, \dots$) in $\phi$.
% \end{restatable}


\begin{restatable}{lemma}{orbitGameSize}\label{lem:orbitGameSize}\label{lem:fnomparsize}
  Given a setup $(\Lcal,\phi_0,s_0,H_0)$ of grade $N$ and $\Gcal_N=\Gcal_N(\Lcal, \phi_0, s_0, H_0)$:
  \begin{enumerate}
    \item the size of $\pos(\gorb(\Gcal_\S))$ is bounded by
    $
        |\orb(\Scal)| \cdot|\phi_0|\cdot \frac{M({\phi_0})!}{|\supp(\phi_0)|!}\cdot (M(\phi_0)+\regindex(\Scal) + 1)^{M(\phi_0) + 1}\cdot(1+o(1))
    $.
    \item the time complexity of constructing $\gorb(\Gcal_N)$ is in
$ O(   (|\phi_0| + \Phi_\Lcal(M(\phi_0)) +\regindex(\Scal))
    \cdot
    \max(M(\phi_0)+\regindex(\Scal) + 2, |\Scal|
    )  \cdot 
    (|\orb(\Scal)| \cdot|\phi_0|\cdot \frac{M(\phi_0)!}{|\supp(\phi_0)|!}\cdot (M(\phi_0)+\regindex(\Scal) + 1)^{M(\phi_0) + 1})
    ^2)$.
\end{enumerate}
\end{restatable}
\begin{proof}  We show~1, while~2 is proven in \cref{app:complexityboundedgame}.
    By definition and \cref{lem:freshClosProps}(1), we know that $\pos(\Gcal_N) \subseteq \fstates \times \clos^*(\phi_0)$.
        Then, by \cref{lem:nomlemma}(1,2) we have that $\pos(\gorb(\Gcal_\S)) \subseteq \orb(\cls{}(\fstates \times \clos^*(\phi_0))) = \orb(\fstates \times \cls{}(\clos^*(\phi_0)))$.
    For every $(s, H, \phi) \in \Gcal_\S$, we have $\supp(s) \subseteq H$ and $|H|\leq N+1$, and so we have $|\orb(\fstates)| \le |\orb(\Scal)| \cdot (\S + 1)$ (what matters of $H$ is the size of $|H\setminus\supp(\phi)|$).
    Using \cref{lem:nomlemma}(3) and the fact that the largest support sizes in $\fstates$ and $\clos^*(\phi_0)$ are bounded by $N+1$ and $|\supp(\phi_0)| + \|\phi_0\|$ (by \cref{lem:freshClosProps}(3)) respectively:
    \begin{align*}
        |\orb(\cls{}(\fstates \times \clos^*(\phi_0))| &\le |\orb(\Scal)| \cdot (\S + 1) \cdot|\orb(\cls{}(\clos^*(\phi_0)))|\cdot (\S + 1)^{|\supp(\phi_0)| + \|\phi_0\|}\cdot(1 + \epsilon) 
        \\ &\le |\orb(\Scal)| \cdot|\orb(\clos^*(\phi_0))|\cdot (\S + 1)^{|\supp(\phi_0)| + \|\phi_0\| + 1}\cdot(1+\epsilon)
\\
        &\le |\orb(\Scal)| \cdot|\phi_0|\cdot \frac{(|\supp(\phi_0)|+\|\phi_0\|)!}{|\supp(\phi_0)|!}\cdot (\S + 1)^{|\supp(\phi_0)| + \|\phi_0\| + 1}\cdot(1+\epsilon)
        % \\
%        &\le |\orb(\Scal)| \cdot|\phi_0|\cdot \frac{(|\phi_0|)!}{|\supp(\phi_0)|!}\cdot (\S + 1)^{|\phi_0| + 1}\cdot(1+\epsilon)
    \end{align*}
    with the last two lines using \cref{lem:nomlemma}(1, second part) and \cref{lem:freshClosProps}(2).
%     and \cref{lem:sizebound}.
\end{proof}

% \begin{restatable}{lemma}{FRAGameSize}\label{app:FRAGameSize}
% 	Let $\Acal$ be an $r$-FRA with fresh-orbit-finite F-LTS $\Lcal$.
% 	Given a setup $(\Lcal,\phi_0,s_0,H_0)$ of grade $N$ and $\Gcal=\Gcal_N(\Lcal, \phi_0, s_0, H_0)$,
% 	the time complexity of constructing $\gorb(\Gcal)$ is in
% 	\begin{align*}
% 		&O(
% 		(2r + |\phi_0| + |\supp(\phi_0)| + \|\phi_0\|)
% 		\cdot
% 		\max(|\supp(\phi_0)| + \|\phi_0\| + r + 2, |Q|
% 		)
% 		\\
% 		&\ \cdot 
% 		(|Q| \cdot|\phi_0|\cdot \frac{(|\supp(\phi_0)|+\|\phi_0\|)!}{|\supp(\phi_0)|!}\cdot (\S + 1)^{|\supp(\phi_0)| + \|\phi_0\| + 1}\cdot(1+o(1)))
% 		^2)
% 	\end{align*}
% \end{restatable}


% \begin{restatable}{lemma}{FRAGameSize}\label{lem:FRAGameSize}
% 	Let $\Acal$ be an $r$-FRA with fresh-orbit-finite F-LTS $\Lcal$.
% 	Given a setup $(\Lcal,\phi_0,(q_0, \rho_0),H_0)$ of grade $N$ and $\Gcal_N=\Gcal_N(\Lcal, \phi_0, (q_0, \rho_0), H_0)$
% 	 the time complexity of constructing $\gorb(\Gcal_N)$ is in 
% 	 $
% 	 O(
% 	 (2r + 2|\phi_0|)
% 	 \cdot
% 	 \max(|\phi_0| + r + 2, |\fstates|
% 	 )
% 	 $
% \end{restatable}
\begin{restatable}{theorem}{MCinFHML}\label{thm:MCinFHML}
  Let $\Lcal$ be a fresh-orbit-finite FLTS with $n$ orbits and set of states $\Scal$, with register index $r$,
% and $(s_0,H_0)\in\hat\Scal$, 
  and let $\phi_0$ be a \fhml\ closed formula with maximum alternation depth $d$. 
    Then, the time complexity of model checking $\phi_0$ on some configuration in $\Lcal$ is in
	\begin{align*}
		&O(\ (n \cdot|\phi_0|\cdot \frac{M(\phi_0)!}{|\supp(\phi_0)|!}\cdot (M(\phi_0)+r + 1)^{M(\phi_0) + 1})^{\log(d) + 6}
		\\ &\ + 
		(|\phi_0| + \Phi_\Lcal(M(\phi_0)) +r)
		\cdot
		\max(M(\phi_0)+r + 2, n
		)\\&\quad \cdot 
		(n \cdot|\phi_0|\cdot \frac{M(\phi_0)!}{|\supp(\phi_0)|!}\cdot (M(\phi_0)+r + 1)^{M(\phi_0) + 1})
		^2\
		)
	\end{align*}
\end{restatable}

\begin{remark}
In the case of model checking formulas on an $r$-FRA $\Acal=\langle Q,\mu,\Sigma,\delta\rangle$, the calculations above are adapted by setting:
\begin{itemize}
\item $|\orb(\Scal)|=|Q|=n$, as two state-assignment pairs are equal up to permutation iff they share the same state;
\item $\regindex(\Lcal)\leq r$, which is equality if there is a state using all $r$  registers;
  \item we can take $\Phi_\Lcal(k)=k+r$, as it suffices to find a matching between the two register assignments (of $s_1$ and $s_2$) and check it complies with mapping $\vec a$ to $\vec b$.
  \end{itemize}
\end{remark}


% \begin{corollary}
%     Given an $r$-FRA $\Acal$ with fresh-orbit-finite F-LTS $\Lcal$, and \fhml\ formula $\phi_0$, 
%     % \fixednote{fresh-orbit-finite?}
%     the time complexity of model checking $\phi_0$ on $\Lcal$ is in
%     \begin{align*}
% 	&O((|Q| \cdot|\phi_0|\cdot \frac{(|\phi_0|)!}{|\supp(\phi_0)|!}\cdot (\S + 1)^{|\phi_0| + 1})^{\log(d) + 6}
% 	%        \\
% 	%        &\ + (
% 	%        (2r + 2|\phi_0|)
% 	%        \cdot
% 	%        \max(|\phi_0| + r + 2, |\fstates|
% 	%        )
% 	\\&\ + (2r + 2|\phi_0|)
% 	\cdot
% 	\max(|\phi_0| + r + 2, |Q|
% 	)
% 	\\
% 	&\ \cdot 
% 	(|Q| \cdot|\phi_0|\cdot \frac{(|\phi_0|)!}{|\supp(\phi_0)|!}\cdot (\S + 1)^{|\phi_0| + 1})
% 	^2))
% \end{align*}
% \end{corollary}

% \begin{proof}
% 	The proof is symmetric to that of the previous theorem. 
% 	From Lemma~\ref{app:FRAGameSize}, the time complexity of constructing $\gorb(\Gcal_\S)$ is in 
% 		\begin{align*}
% 		&O(
% 		(2r + |\phi_0| + |\supp(\phi_0)| + \|\phi_0\|)
% 		\cdot
% 		\max(|\supp(\phi_0)| + \|\phi_0\| + r + 2, |Q|
% 		)
% 		\\
% 		&\ \cdot 
% 		(|Q| \cdot|\phi_0|\cdot \frac{(|\supp(\phi_0)|+\|\phi_0\|)!}{|\supp(\phi_0)|!}\cdot (\S + 1)^{|\supp(\phi_0)| + \|\phi_0\| + 1}\cdot(1+o(1)))
% 		^2)
% 	\end{align*}
%     The number of orbits in $\Lcal$ is $|Q|\cdot r \le |Q| \cdot (\S + 1)$, and so it follows from the previous theorem and Lemma~\ref{lem:sizebound} that the final complexity is in:
%     \begin{align*}
%         &O((|Q| \cdot|\phi_0|\cdot \frac{(|\phi_0|)!}{|\supp(\phi_0)|!}\cdot (\S + 1)^{|\phi_0| + 1})^{\log(d) + 6}
% %        \\
% %        &\ + (
% %        (2r + 2|\phi_0|)
% %        \cdot
% %        \max(|\phi_0| + r + 2, |\fstates|
% %        )
%         \\&\ + (2r + 2|\phi_0|)
%         \cdot
%         \max(|\phi_0| + r + 2, |Q|
%         )
%         \\
%         &\ \cdot 
%         (|Q| \cdot|\phi_0|\cdot \frac{(|\phi_0|)!}{|\supp(\phi_0)|!}\cdot (\S + 1)^{|\phi_0| + 1})
%         ^2))
%     \end{align*}
% \end{proof}

% <<<<<<< HEAD
% =======
% \begin{proof}
% 	The proof is symmetric to that of the previous theorem. 
% 	From Lemma~\ref{app:FRAGameSize}, the time complexity of constructing $\gorb(\Gcal_\S)$ is in 
% 		\begin{align*}
% 		&O(
% 		(2r + |\phi_0| + |\supp(\phi_0)| + \|\phi_0\|)
% 		\cdot
% 		\max(|\supp(\phi_0)| + \|\phi_0\| + r + 2, |Q|
% 		)
% 		\\
% 		&\ \cdot 
% 		(|Q| \cdot|\phi_0|\cdot \frac{(|\supp(\phi_0)|+\|\phi_0\|)!}{|\supp(\phi_0)|!}\cdot (\S + 1)^{|\supp(\phi_0)| + \|\phi_0\| + 1})
% 		^2)
% 	\end{align*}
%     The number of orbits in $\Lcal$ is $|Q|\cdot r \le |Q| \cdot (\S + 1)$, and so it follows from the previous theorem and Lemma~\ref{lem:sizebound} that the final complexity is in:
%     \begin{align*}
%         &O((|Q| \cdot|\phi_0|\cdot \frac{(|\phi_0|)!}{|\supp(\phi_0)|!}\cdot (\S + 1)^{|\phi_0| + 1})^{\log(d) + 6}
% %        \\
% %        &\ + (
% %        (2r + 2|\phi_0|)
% %        \cdot
% %        \max(|\phi_0| + r + 2, |\fstates|
% %        )
%         \\&\ + (2r + 2|\phi_0|)
%         \cdot
%         \max(|\phi_0| + r + 2, |Q|
%         )
%         \\
%         &\ \cdot 
%         (|Q| \cdot|\phi_0|\cdot \frac{(|\phi_0|)!}{|\supp(\phi_0)|!}\cdot (\S + 1)^{|\phi_0| + 1})
%         ^2))
%     \end{align*}
% \end{proof}
% >>>>>>> refs/remotes/origin/main

%%% Local Variables:
%%% mode: latex
%%% TeX-master: "main"
%%% End:
